图像识别是人工智能和计算机视觉领域的重要研究方向，在智慧安防、医学辅助诊断、工业质检和无人驾驶等场景中具有广泛的应用价值。随着深度学习的发展，卷积神经网络能够自动完成多层次特征提取，在公开数据集上的识别性能显著优于传统手工特征方法。然而，现有模型仍然存在参数规模大、部署成本高以及复杂场景下泛化能力不足等问题，因此针对识别精度与模型效率的协同优化开展研究具有现实意义。

本文围绕“轻量化图像识别模型设计”展开研究，首先梳理卷积神经网络、残差学习和注意力机制的基本原理，对经典图像识别方法的技术路线、适用场景与局限性进行归纳分析；其次，在基准卷积网络基础上构建一种融合通道重标定与多尺度特征聚合的改进模型，通过引入深浅层信息交互机制增强目标区域表征能力；再次，结合数据增强、迁移学习和消融实验设计，建立完整的训练与评估流程，并从准确率、参数量、推理时延等角度对模型性能进行量化比较。

实验结果表明，所提出方法在公开分类数据集上的识别精度较基线网络有明显提升，同时保持了较低的参数规模与推理开销。通过对错误样本的可视化分析可以发现，多尺度信息融合策略能够有效缓解复杂背景、尺度变化和局部遮挡对模型判别能力的影响。本文研究说明，在本科毕业设计规模下，围绕模型结构、训练策略与实验分析构建一套完整的方法链路，能够较好体现毕业设计（论文）对研究过程、实验论证和规范表达的综合要求。

\gxukeywords{图像识别；深度学习；卷积神经网络；轻量化模型；注意力机制}
